\clearpage
\thispagestyle{empty}
\pdfbookmark[0]{Instructions for Determining the Duration of a Magnetic Needle’s Swing. Edition and sources}{magnetic-needle-swing-sources}
\begin{center}\Large\textbf{Instructions for Determining the Duration of a Magnetic Needle’s Swing}\end{center}

Carl Friedrich Gauss’s instructions appeared in the \textit{Resultate aus den Beobachtungen des magnetischen Vereins} for 1837, article IV. This AI translation and AI-produced typesetting follow the complete German text in \textit{Werke}, vol.~V (Göttingen, 1867), the unnumbered opening and pp.~375--394, including the five historical notes. A swing here runs from one extreme position to the next. The observations and formulas retain their historical notation and units.

\smallskip\noindent
For terminology on passage instruments and timing, see the Deutsches Museum catalogue entry \textit{Passageinstrument}, edited by Julia Bloemer and Benjamin Mirwald, no.~701, version of 24 August 2016, and its glossary entries \textit{Auge-Ohr-Methode} and \textit{Fadenkreuz}. For modern German terminology on damped oscillations, see the Karlsruhe Institute of Technology physics laboratory manuals linked below.

\medskip\noindent
\href{https://gdz.sub.uni-goettingen.de/id/PPN236006339}{Collected German edition (1867)}\par
\noindent\href{https://digital.deutsches-museum.de/projekte/gruendungssammlung/detail/701/}{Passage instrument catalogue entry}\par
\noindent\href{https://digital.deutsches-museum.de/projekte/gruendungssammlung/glossar/}{Deutsches Museum glossary}\par
\noindent\href{https://labs.physik.kit.edu/prakt-klass-physik-manuals/P1/Resonanz/doc/Hinweise-Schwingung/}{KIT Classical Physics Laboratory, notes on oscillations}\par
\noindent\href{https://labs.physik.kit.edu/prakt-nebenfach-manuals/Experiments/Erzwungene_Schwingungen/}{KIT Physics Laboratory, forced oscillations}
\clearpage
